\section{The linear equation and the firing trace}\label{sec:linear}

For the construction we fix
\begin{equation}\label{eq:physical-choice}
 v_E=2,\qquad v_F=1,\qquad y_L=1.
\end{equation}
Thus $D=(0,1)\times(0,\infty)$,
$J(v,y)=y(2-v)-v$, $y_F=1$, $J_0(y)=2y$ and $J_F(y)=y-1$.
We first keep the conductance mean and diffusion prescribed. Write them as
$\mu(t)$ and $\mathfrak a(t)$, bounded above and below by positive constants.
The conductance flux is
\[
 Q_{\mu,\mathfrak a}f=(\mu-y)f-\mathfrak a f_y.
\]
The weak formulation we use retains this flux at the reflecting boundary.
If $\gamma_0 f,\gamma_F f$ are the voltage traces, they satisfy
\begin{equation}\label{eq:voltage-traces}
 \gamma_0 f=\gamma_F f=0\quad(0<y<1),\qquad
 J_0\gamma_0 f=J_F\gamma_F f\quad(y>1).
\end{equation}
For a smooth test $\varphi$, compactly supported in $y$ and allowed to be
nonzero at all physical faces, the equation reads
\begin{align}\label{eq:weak-prescribed}
 \left[\int_D f\varphi\right]_{t_1}^{t_2}
 &+\int_{t_1}^{t_2}\!\int_1^\infty
 J_F\gamma_F f\,[\varphi(t,1,y)-\varphi(t,0,y)]\dd y\dd t\notag\\
 &=\int_{t_1}^{t_2}\!\int_D
 f\,[\varphi_t+J\varphi_v+(\mu-y)\varphi_y]
       -\mathfrak a f_y\varphi_y.
\end{align}
There is no boundary term at $y=0$. This is the zero conormal-flux condition.
For periodic solutions the endpoint term cancels when the test is periodic.
All voltage conditions are understood almost everywhere in their kinetic
trace measures. No trace value is needed at $y=1$.

\subsection{The frozen evolution}

Let $m$ belong to $I=[2,3]$ and let $a>0$. Fix $\sigma>1$ with
$2a\sigma<m$; we eventually take $\sigma=2$. Define
\[
 W(y)=e^{2\sigma y},\qquad
 \HH_\sigma=L^2(D,W\dd v\dd y),\qquad
 \HH_{-\sigma}=L^2(D,W^{-1}\dd v\dd y).
\]
\[
 \HH_{\sigma,0}=\left\{h\in\HH_\sigma:\int_Dh=0\right\}.
\]
We complexify these spaces. Their transpose pairing is the bilinear integral
$\langle h,u\rangle=\int_Dhu$; conjugation is used only in Hilbert energy
identities. The mass functional is bounded on $\HH_\sigma$.

We use the following consequences of the prescribed linear theory of
Fonte Sanchez and Mischler.

\begin{proposition}[The frozen linear evolution]\label{prop:frozen}
For prescribed positive bounded $\mu,\mathfrak a$ and initial data in
$\HH_\sigma$, there is a unique solution of the linear equation in
$C([0,t_*];\HH_\sigma)$, with local $L^2$ conductance derivative and the
renormalized weak traces, for every finite $t_*$. It preserves positivity
and mass. For $\mu=m$, $\mathfrak a=a$, its evolution is a strongly continuous
semigroup $S(t)$ on $\HH_\sigma$ with generator $\LL=\LL_{m,a}$.
There is a normalized nonnegative invariant density $M=M_{m,a}$, belonging
to every fixed exponential weighted space, and constants $C_a,\kappa_a>0$
such that
\begin{equation}\label{eq:frozen-gap}
 \norm{S(t)h}_{\HH_\sigma}\le C_a e^{-\kappa_a t}
           \norm{h}_{\HH_\sigma},\qquad h\in\HH_{\sigma,0}.
\end{equation}
The stationary density is in $D(\LL)$ and $\LL M=0$.
\end{proposition}

\begin{proof}
The prescribed statement and its renormalized Green identity are
\cite[Theorem 2.2 and (2.18)--(2.19)]{FM24}. The stationary and decay
statements follow from \cite[Theorem 1.2]{FM24} by setting the source's
connectivity to zero and its background mean and diffusion to $m,a$.
In particular its coefficients are then exactly $m-y,a$. The invariant
state is the generator-domain stationary state constructed in Section 7
of that paper. Autonomy and uniqueness give the semigroup property;
the asserted time continuity gives strong continuity. Its $p=2$ decay
estimate is \eqref{eq:frozen-gap}. These are estimates for the frozen
equation. The firing feedback will be included in Section~\ref{sec:crossing}.
\end{proof}

The domain $D(\LL)$ means the strong semigroup generator domain. We do not
replace it by a prescribed collection of smooth functions at the two
characteristic corners. The differential expressions in the interior are
\[
 \LL h=-\partial_v(Jh)-\partial_y((m-y)h)+ah_{yy},
 \qquad
 \G u=Ju_v+(m-y)u_y+au_{yy}.
\]

The outgoing space and the firing functional are
\begin{equation}\label{eq:out-space}
 \OO_\sigma=L^2((1,\infty),J_FW\dd y),\qquad
 \ell h=\int_1^\infty J_F(y)h_F(y)\dd y.
\end{equation}
The next estimate defines these quantities on the actual generator domain.

\begin{lemma}[Forward energy and the generator trace]\label{lem:forward-energy}
There is a constant $C=C_{m,a,\sigma}$ such that
\begin{align}\label{eq:forward-energy}
 \norm{S(t)h}_{\HH_\sigma}^2
 &+2a\int_0^t\norm{\partial_yS(s)h}_{\HH_\sigma}^2\dd s
 +\int_0^t\!\int_1^\infty J_F(1-J_F/J_0)
       |(S(s)h)_F|^2W\dd y\dd s\notag\\
 &\le \norm h_{\HH_\sigma}^2+
 C\int_0^t\norm{S(s)h}_{\HH_\sigma}^2\dd s.
\end{align}
For $h\in D(\LL)$, its derivative and outgoing trace satisfy
\begin{equation}\label{eq:domain-energy}
 2a\norm{h_y}_{\HH_\sigma}^2+\tfrac12\norm{h_F}_{\OO_\sigma}^2
 \le C\norm h_{\HH_\sigma}^2-2\operatorname{Re}(\LL h,h)_{\HH_\sigma}.
\end{equation}
In particular $\ell$ is continuous for the graph norm of $D(\LL)$.
The finite-time version of \eqref{eq:forward-energy} is uniform for
prescribed coefficients in a sufficiently small positive neighborhood of $m,a$.
\end{lemma}

\begin{proof}
Here are the truncations in the Green identity. For real $f$ use
\[
 b_n(s)=2n^2\bigl(\sqrt{1+s^2/n^2}-1\bigr),\qquad
 D_n(s)=s b_n'(s)-b_n(s).
\]
These are smooth convex functions with bounded first and second
derivatives, $0\le D_n\le b_n\le s^2$, and $b_n''\to2$. If boundedness of
the renormalization itself is required, first compose $b_n$ with a smooth
bounded argument truncation at height $A$. For fixed $n$ its first and
second derivatives are bounded independently of $A$. On compact physical
intervals the local $L^2$ derivative and the source trace bounds dominate
every term, so $A\to\infty$ recovers $b_n$. This also applies to the
entropy renormalizations used later.

Use $b_n$ with $W$ cut off between $R$ and $2R$. At high conductance the
reset condition, with $r=J_F/J_0\in(0,1)$, makes the boundary loss
\[
 J_F\bigl[b_n(f_F)-r^{-1}b_n(rf_F)\bigr]\ge0.
\]
It converges to $J_F(1-r)|f_F|^2$. Integrate the cross diffusion term as
$-a\int\partial_y b_n(f)W_y$. At the floor this gives
$-mD_n(f_0)+2a\sigma b_n(f_0)$. The source local $H^1_y$ class gives
$f_0\in L^2(\dd t\dd v)$, so this term converges to
$-(m-2a\sigma)|f_0|^2$. Its sign is used after, rather than before,
the truncation is removed.

Away from the cutoff the volume term is
\[
 (y+2)D_n(f)+2\sigma(m-y)b_n(f)+4a\sigma^2 b_n(f)
 \le C|f|^2,
\]
since $2\sigma>1$. The $y$-cutoff errors are bounded by a weighted
$\HH_\sigma$ tail: on their support $y/R$ is bounded. Pass $n\to\infty$
for fixed $R$, then $R\to\infty$. Fatou's lemma retains the nonnegative
gradient and reset losses. This proves \eqref{eq:forward-energy} and also
the favorable floor loss. Apply the argument to real and imaginary parts.

For $h\in D(\LL)$ divide by $t$. The norm difference quotient tends to
$2\operatorname{Re}(\LL h,h)$. The averages
$t^{-1}\int_0^t S(s)h\dd s$ tend to $h$ in graph norm. Jensen's inequality
and the preceding estimate bound their gradients and traces weakly;
their averaged Green identities identify the limits. As
$1-J_F/J_0\ge1/2$, weak lower semicontinuity gives
\eqref{eq:domain-energy}. Finally
\[
 |\ell h|\le
 \left(\int_1^\infty J_FW^{-1}\dd y\right)^{1/2}\norm{h_F}_{\OO_\sigma}.
\]
All coefficients and the limiting floor sign remain controlled in a
small positive prescribed-coefficient neighborhood.
\end{proof}

\subsection{Divergence inputs and the reflecting boundary}

The linearized feedback is a conductance divergence. For compact smooth
$g$ supported inside $D$, put $Bg=\partial_y g$. The completion for
general $g\in\HH_\sigma$ will be defined by
\begin{equation}\label{eq:B-definition}
 \langle Bg,\varphi\rangle=-\int_Dg\varphi_y,
\end{equation}
including tests nonzero at the floor. In particular it does not assign
a point value to $g(0)$.

By \eqref{eq:frozen-gap}, on the zero-mass space the ordinary inverse is
\begin{equation}\label{eq:ordinary-resolvent}
 R_{i\theta}h=\int_0^\infty e^{-i\theta t}S(t)h\dd t,
 \qquad \theta\in\R.
\end{equation}
It maps into $D(\LL)$ and $(i\theta-\LL)R_{i\theta}h=h$.
This follows by taking the difference quotient of the integral under
$S(s)$; the omitted integral over $(0,s)$ tends to $h$ after division by
$s$. The gap also excludes a homogeneous zero-mass solution.

We first bound divergence inputs independently of a killed backward
harmonic. The bounded-source backward construction needed in this proof
is given in Appendix~\ref{app:boundary}; in particular it proves the
actual Neumann condition before taking a global dual energy estimate.

\begin{proposition}[Completed divergence resolvent]\label{prop:divergence}
There is a bounded map
\[
 \mathcal R_\theta:g\longmapsto(h,h_y,h_F)
 \quad\hbox{from }\HH_\sigma\hbox{ to }
 \EE_\sigma:=\HH_{\sigma,0}\times\HH_\sigma\times\OO_\sigma,
\]
uniformly for all $\theta\in\R$, whose density component extends
$h=R_{i\theta}Bg$ for compact smooth $g$. Its weak equation is
\begin{align}\label{eq:completed-resolvent}
 i\theta\int_Dh\varphi
 -\int_Dh[J\varphi_v+(m-y)\varphi_y]
 +a\int_Dh_y\varphi_y
 &+\int_1^\infty J_Fh_F(\varphi_F-\varphi_0)\dd y\notag\\
 &=-\int_Dg\varphi_y.
\end{align}
Thus the natural flux $(m-y)h-ah_y-g$ has zero floor trace.
The original voltage conditions hold in the completed trace sense.
We continue to denote the density component by $R_{i\theta}Bg$.
The maps extend holomorphically to a neighborhood of each fixed
imaginary frequency where the ordinary zero-mass resolvent exists.
Constants and neighborhood sizes may depend on the fixed $a>0$.
\end{proposition}

\begin{proof}
Let $s$ exceed the full semigroup growth bound and a fixed dual-energy
constant. Here $Z_t=(V_t,Y_t)$ has reflected conductance
$\dd Y_t=(m-Y_t)\dd t+\sqrt{2a}\dd W_t+\dd K_t$ and voltage
$\dot V_t=J(V_t,Y_t)$; a threshold hit resets voltage from $1$ to $0$
without changing conductance. Proposition~\ref{prop:process} establishes
its density law, and Proposition~\ref{prop:backward-source} supplies
the transpose-resolvent representation below.
For bounded compact smooth $r$ the full-process value
\[
 u(z)=\E_z\int_0^\infty e^{-(s+i\theta)t}r(Z_t)\dd t
\]
is the transpose of the full forward resolvent. It is bounded by
$\norm r_\infty/s$. Proposition~\ref{prop:backward-source} in the appendix
gives, by localization at the corners and the actual Neumann floor,
\begin{equation}\label{eq:large-real-dual}
 (s-C)\norm u_{\HH_{-\sigma}}^2+
 a\norm{u_y}_{\HH_{-\sigma}}^2
 \le\norm r_{\HH_{-\sigma}}\norm u_{\HH_{-\sigma}}.
\end{equation}
Completion of compact smooth $r$, using closedness of the distributional
derivative, gives the same estimate for every $r\in\HH_{-\sigma}$.

For an arbitrary such $r$ set $r_0=r-(\int_DMr)\one$. The transpose of
\eqref{eq:ordinary-resolvent}, represented in $\int_D M\psi=0$, gives
\[
 (i\theta-\LL^T)\psi=r_0,\qquad
 \norm\psi_{\HH_{-\sigma}}\le C_a\norm r_{\HH_{-\sigma}}.
\]
The representative is an actual transpose-domain element: decompose
$h\in D(\LL)$ into $(\int h)M+h_0$. The stationary terms vanish by the
two mean-zero normalizations. The shifted identity is
\[
 \psi=(R^{\rm full}_{s+i\theta})^T(r_0+s\psi).
\]
Its input is bounded uniformly in $\theta$. Applying
\eqref{eq:large-real-dual} at this same frequency proves
$\norm{\psi_y}_{\HH_{-\sigma}}\le C_a\norm r_{\HH_{-\sigma}}$.

For compact $g$, $h=R_{i\theta}Bg\in D(\LL)$ and transpose duality now
gives
\[
 \langle h,r\rangle=\langle Bg,\psi\rangle=-\int_Dg\psi_y.
\]
Taking the dual norm bounds $\norm h_{\HH_\sigma}$ by $C_a\norm g_{\HH_\sigma}$.
Use \eqref{eq:domain-energy} with
$\LL h=i\theta h-Bg$. The remaining source term is
\[
 (Bg,h)_{\HH_\sigma}
   =-\int_Dg(\overline{h_y}+2\sigma\overline h)W.
\]
Young's inequality absorbs its derivative term. Hence
\begin{equation}\label{eq:triple-bound}
 \norm h_{\HH_\sigma}+\norm{h_y}_{\HH_\sigma}
       +\norm{h_F}_{\OO_\sigma}\le C_a\norm g_{\HH_\sigma},
\end{equation}
uniformly in $\theta$, including zero.

Approximate general $g$ by compact smooth functions away from $y=0$.
Apply \eqref{eq:triple-bound} to differences. All three components
converge strongly, and their weak equations give
\eqref{eq:completed-resolvent}. This is why the inhomogeneous floor flux
is retained even though the approximants vanish there. The product
$\EE_\sigma$ is an ambient Hilbert space; derivative and trace identities
are consequences of the weak equation for these outputs.

Finally, for $\lambda_0=i\theta$ use
\[
 R_\lambda B=R_{\lambda_0}B+
       (\lambda_0-\lambda)R_\lambda R_{\lambda_0}B.
\]
The second term has an ordinary $\HH_{\sigma,0}$ datum, so the ordinary
resolvent identity and domain energy give the local holomorphic extension.
\end{proof}

We will also need a finite-time form of the same completion.

\begin{lemma}[Causal divergence forcing]\label{lem:causal}
For $g\in L^2((0,t_*);\HH_\sigma)$ there is a unique zero-initial-data
solution of $(\partial_t-\LL)w=Bg$ in the completed energy class, with
\begin{equation}\label{eq:causal-bound}
 \sup_{0\le t\le t_*}\norm{w(t)}_{\HH_\sigma}^2+
 \int_0^{t_*}\bigl(\norm{w_y}_{\HH_\sigma}^2+
                  \norm{w_F}_{\OO_\sigma}^2\bigr)\dd t
 \le C_{a,m,t_*}\norm g_{L^2_t\HH_\sigma}^2.
\end{equation}
Its mass is zero and its natural base flux equals $g$ at the floor in
the weak formulation.
\end{lemma}

\begin{proof}
First take $g$ continuously differentiable in time and compactly supported
inside $D$. Its divergence is a differentiable $\HH_\sigma$ forcing, so
the semigroup convolution is a classical $D(\LL)$ solution. This also
follows by integrating the time derivative of $Bg$ in that convolution.
Apply \eqref{eq:domain-energy} to $w(t)$. Integrating the input as above,
Young's inequality and Gronwall give \eqref{eq:causal-bound}.
Complete the inputs in $L^2_t\HH_\sigma$. The outputs converge in
$C_t\HH_\sigma$, and their derivatives and traces converge in $L^2$.
Their weak equations retain the same floor input. The difference of two
completions is a homogeneous zero-initial-data solution, so prescribed
uniqueness in Proposition~\ref{prop:frozen} proves uniqueness. All these
solutions lie in its local energy class; a smaller exponential weight
may be used for this uniqueness comparison.
\end{proof}
